% !TEX root = ../main.tex

\chapter{Background: Data, Compute, and Feedback Bottlenecks in BSP Systems}
\label{chap:bg}
\minitoc

We first review key concepts and challenges in the bulk-synchronous parallel (BSP) paradigm,
beginning with an overview of the computing environment (\S\ref{sec:bg:hpc}), followed by
data management challenges (\S\ref{sec:bg:data}), compute efficiency challenges
(\S\ref{sec:bg:compute}), and limitations of existing feedback
mechanisms (\S\ref{sec:bg:feedback}).
While primarily focusing on scientific computing, we also discuss distributed machine learning
training workloads, as their shared bulk-synchronous characteristics create similar challenges and
inform our later work.

% We first review key concepts and challenges in the bulk-synchronous parallel (BSP) paradigm,
% beginning with an overview of the computing environment (\S\ref{sec:bg:hpc}), followed by a review
% of data management (\S\ref{sec:bg:data}) and compute efficiency challenges
% (\S\ref{sec:bg:compute}).
% % Maybe note ML differences as a separate section
% While we primarily focus on scientific computing, we also discuss distributed machine learning
% training workloads, as their shared bulk-synchronous characteristics create similar challenges and
% inform our later work.

\section{BSP: Execution, Architecture, and Applications}
\label{sec:bg:hpc}

\todo[inline]{Move BSP up}
\todo[inline]{In general make this better}

\begin{figure}[t]
  \centering
  \includegraphics[width=1.0\textwidth]{bsp-arch.pdf}
  \caption{BSP architecture and execution model.
    (Left) Applications consist of CPU or GPU ranks, connected via an RDMA fabric and parallel storage.
    (Right) Execution proceeds in timesteps followed by global synchronization via collectives.
    Runtime is therefore dictated by the slowest rank.
    Periodically, checkpoints are triggered for fault tolerance and offline analysis.
  }
  \label{fig:bg:arch}
\end{figure}

The BSP model consists of alternating compute and synchronization phases (see \cref{fig:bg:arch}).
% \paragraph{Execution Model}
Execution begins with decomposition of the problem into independent subproblems assigned to
workers, and proceeds as discrete timesteps.
Each timestep consists of a compute phase followed by a global synchronization phase, where
synchronous collectives~\cite{:MPI, :NCCL} ensure all ranks advance together.
Every few timesteps, a checkpoint is triggered for both fault tolerance and offline analysis.
This tight coupling is the root of cluster efficiency problems: synchronous checkpoints create
coordinated I/O bursts that stress the storage system, while global collectives make overall
performance acutely sensitive to the slowest worker.
% As clusters scale, checkpointing frequency increases as failure probability rises 
Growth in cluster sizes leads to increased checkpointing frequency, as the increased failure
probabilty makes forward progress more valuable~\cite{Jones2012:Checkpointing}.
% Checkpointing frequency increases with application scale, intra-checkpoint failure probability
% rises, creating periodic I/O-bound operations that stress the storage system. Checkpointing
% frequency increases with application scale as intra-checkpoint failure probability rises, creating
% periodic I/O-bound operations that stress the storage system.


\paragraph{HPC Cluster Architecture and Applications}
Large-scale HPC clusters typically consist of O(10--100K) nodes~\cite{:TOP500ListJune} connected by
a
high-performance fabric~\cite{Pfist2001:Infiniband,Rowet2022:Slingshot,Metz2024:UltraEthernet}.
Each node typically consists of CPUs and typically multiple GPUs, which are increasingly important
for scientific computing~\cite{DeBar2014:HPCGPU,Snell2017:HPCGPU}.
Storage is traditionally provided by parallel filesystems~\cite{Braam2002:Lustre} but object
stores~\cite{Henne2020:DAOS} and richer commercial offerings~\cite{2024:Hammerspace,:VASTPlatform}
are increasingly common~\cite{Glenn2025:NoPFS}.

% \paragraph{Applications}
Scientific applications target fundamental questions across domains such as cosmology and climate
science, problems that require simulating complex physical phenomena over extended spatial and
temporal scales.
These applications can occupy the majority of a supercomputer's resources for months at a
time~\cite{Das2023:LargeScale, Stock2024:LargeScale,Gary2023:LANLPlatformPlanning}.
From a systems perspective, they can be simplified into two common archetypes: \emph{particle
  codes}~\cite{Bird2022:VPIC2.0} and \emph{mesh codes}~\cite{Dubey2014:Survey}.
Particle codes track the state and interactions of billions individual discrete entities as they
move through the simulation domain.
Mesh codes discretize the spatial domain into grids and compute the evolution of properties of
interest at each grid point.


\paragraph{Distributed ML Training}
Distributed ML training represents a similar and increasingly important cluster workload, occupying
tens of thousands of GPUs and running for
months~\cite{:LlamaScale,Rahma2024:MLScale,Gratt2024:Llama3Scale}.
The cluster architectures are comparable, though ML clusters often add a \emph{frontend} Ethernet
network for general datacenter connectivity~\cite{Gangi2024:RDMA}.
While the parallelization strategies differ---ML employs a mix of data, tensor, and pipeline
parallelism---the execution model still relies on the same core BSP primitives of synchronous
collectives~\cite{Gangi2024:RDMA} and periodic checkpointing~\cite{Eisen2022:Checkpoint}.

\section{Data Management Challenges}
\label{sec:bg:data}

Checkpointing creates I/O-bound bottlenecks in BSP systems, as coordinated writes across thousands
of ranks stress parallel filesystems and consume significant
runtime~\cite{Xie2012:CharacterizingOutputBottlenecks}.
Traditional workflows then require scientists to re-read and analyze this data after simulation
completion, consuming additional I/O bandwidth and computational resources while delaying
scientific discovery.
Ingesting data using conventional databases is not viable, as databases incur multiple on-disk
writes for each application write~\cite{Raju2017:PebblesDB} (\emph{write amplification}), offering
poor tradeoffs for ingestion-heavy workloads.

In-situ processing~\cite{Child2020:TerminologySituVisualization} analyzes data while it streams
from application to storage, reducing both I/O overhead and time-to-insight.
This approach is employed for multiple scientific use cases ranging from
visualization~\cite{Ayach2015:Catalyst,Ayach2016:SENSEI,Larse2017:ALPINE,Larse2022:Ascent} to
indexing~\cite{Eisen2024:ADIOS,Docan2010:DataSpaces}.
In-situ reorganization using hash-partitioning and shuffling has been demonstrated to be viable at
HPC scale~\cite{Zheng2018:DeltaFS}.
However, unlike databases that provide general query interfaces over composable indexes, existing
in-situ systems are purpose-built for specific use cases and lack programmability.
A critical gap remains in range query support, essential for analytics on continuous physical
attributes but absent from existing in-situ frameworks.

% Checkpointing creates I/O-bound bottlenecks in BSP systems. Traditional post-processing analytics
% requires additional time and computational resources after simulation completion. Databases, while
% providing composable interfaces, suffer from high write amplification factors (WAF) that make them
% unsuitable for HPC ingestion workloads. This motivated the development of in-situ processing
% approaches.

% In-situ visualization systems (ParaView Catalyst, VisIt) are widely deployed and mature. In-situ
% data reorganization approaches have demonstrated viability, such as DeltaFS hash partitioning at
% large scale. These systems are purpose-built and specialized, creating an ecosystem where
% approaches cannot be easily mixed and matched due to their specialized nature. A critical gap
% remains in range query support, which is essential for spatial and temporal analytics.

% Ingestion efficiency essential. WAF of 1X ideal and essential to ingest at max storage system
% bandwidth. Databases are not ideal.

% Special-purpose systems used for in-situ aggregation and visualization. In-situ indexing and
% reordering less widespread but demonstrated to be viable. Deltafs (insitu hash partitioning). But
% too much reliance on bitmap indices etc, as partitioning is expensive. Partitioning for range
% queries is hard: this currently uses post-processing sorts, which need to read all data again and
% write them, requiring extra hardware and delaying time to insight.

% The biggest missing piece in the in-situ toolking is range partitioning.

\section{Compute Efficiency Challenges}
\label{sec:bg:compute}

Cluster utilization challenges persist in BSP systems: a study of exported traces found more than
60\% of runtime spent in MPI~\cite{:MPI} operations (communication and synchronization), even at
modest scales.
We now provide additional context on performance tooling and workload characterization bottlenecks
that contribute to these challenges.

% \subsection{Performance Analysis Tools}

\paragraph{Performance Analysis Tools}
While real time cluster-level monitoring is common~\cite{Schwa2021:LDMS}, application-level
performance analyses are largely
post-mortem~\cite{Wylie2025:15Years,Shend2006:TAU,2025:kineto,:HTA,Knupf2012:Score-P,Boehm2016:Caliper}.
Two main approaches are tracing and profiling.
Tracing generates fine-grained event logs for each worker (\emph{rank}), which are then aggregated
into a composite view via postprocessing.
The volume of generated trace data presents an analytics challenge, delaying
insight~\cite{Devar2024:DFTracer}, and remains unviable at the largest
scales~\cite{Cui2026:Flare}.
In comparison, profiling generates phase-wise summaries of execution for each rank.
While far more scalable, this obscures fine-grained interactions that degrade tightly-coupled BSP
execution.
% Current tools rely on three main approaches: sampling provides statistical summaries at regular
% intervals, profiling offers detailed function-level breakdowns, and tracing generates
% comprehensive execution records. Sampling and profiling impose arbitrary aggregations, obscuring
% problem structure and temporal patterns necessary for optimization. Tracing produces massive
% datasets in unstructured formats requiring extensive offline processing. This methodology
% inherently operates on historical data, providing no mechanism for real-time performance insights
% or adaptive runtime optimization during execution. Current tools rely on three main approaches:
% sampling provides statistical summaries at regular intervals, profiling offers detailed
% function-level breakdowns, and tracing generates comprehensive execution records. Sampling and
% profiling impose arbitrary aggregations, obscuring problem structure and temporal patterns
% necessary for optimization. Tracing produces massive datasets in unstructured formats requiring
% extensive offline processing. This methodology inherently operates on historical data, providing
% no mechanism for real-time performance insights or adaptive runtime optimization during execution.

\paragraph{Workload Characterization Bottlenecks}
While efficiency challenges are known and attributed to load
imbalance~\cite{Dubey2014:Survey,Burst2008:ALPS,Klenk2017:OverviewMPICharacteristics}, deeper
understanding remains elusive, partly because of visibility bottlenecks discussed above.
Characterization of workload properties such as compute kernel distributions, communication
bottlenecks, placement and load-balancing requirements remain poorly understood.
This is further exacerbated by the separation of analyses from data collection, as the latter are
undertaken by separate groups without access to the underlying hardware that generated the
data~\cite{:DOEProxyApps}.
As a result, distinguishing workload characteristics from environment artifacts becomes even more
difficult.
% As a result of visibility bottlenecks, efficiency challenges, while commonly ascribed to load
% imbalance, While efficiency challenges are commonly ascribed to load imbalance, fundamental
% aspects of this problem remain poorly understood. Application properties, computational
% requirements, and the constraints that must be honored during optimization are not well
% characterized. Fine-grained performance studies typically involve exporting large representative
% traces, which are then analyzed by other research groups often without access to the underlying
% hardware characteristics or runtime context needed for actionable insights. This disconnect
% between data collection and analysis creates fundamental barriers to understanding system
% behavior, leaving the actual nature of performance bottlenecks obscured despite widespread
% acknowledgment of their importance.

% \subsection{Workload Characterization Bottlenecks} Performance analysis in large-scale BSP systems
% relies on aggregate metrics and post-hoc trace analysis. Fine-grained performance artifacts can
% only be captured through comprehensive tracing, as aggregate metrics obscure the detailed
% execution patterns necessary for optimization. Execution snapshots generate substantial trace
% datasets that are typically published for offline analysis by performance analysis groups.
% However, trace analysis is often conducted by researchers without direct access to the underlying
% hardware characteristics or runtime context, creating a disconnect between the data and the
% systems knowledge needed for actionable insights. This methodology inherently operates on
% historical data rather than providing real-time performance insights.

% \subsection{Performance Analysis Tools} Sampling tools provide statistical summaries, profiling
% tools offer detailed breakdowns, and tracing tools generate comprehensive execution records stored
% in formats such as OTF2 for HPC and PyTorch/Kineto traces for ML training. These raw trace formats
% produce large datasets that are not immediately amenable to interactive analysis. Similar
% trace-based methodologies are employed across HPC and ML training domains as scale increases, with
% ML clusters increasingly facing comparable runtime analysis challenges. Performance optimizations
% currently require domain-specific tooling with hardcoded analysis logic, and existing tools do not
% support real-time performance monitoring or adaptive runtime optimization.

% \subsection{Workload Characterization} Reliance on aggregate metrics. Did this improve runtime
% overall?

% Fine-grained effects hard to untangle, but esesential, as our work shows. Only available with
% tracing.

% Trace analysis model: a small snapshot of execution results in a large trace. Published and
% analyzed by other groups. This makes learnings slow and disconnected from hardware. Those who
% analyze traces are not able to distinguish fundamental characteristics from tunable runtime
% artifacts.

% \subsection{Perf Tools}

% Perf tools can be sampling, profiling, or tracing. Former emit summaries. Latter emits lots of raw
% data in formats like OTF2. These are not analytical formats and are not amenable to analytics.

% ML: Pytorch/Kineto traces. Clusters starting to eclipse HPC. Similar runtime challenges.

% Solutions for perf challenges are hardcoded assertions in different tooling.




\section{Feedback Mechanisms and Limitations}\label{sec:bg:feedback}


\paragraph{Computational Steering}
Computational steering refers to interactive workflows where output from a running simulation
is observed and parameters are adjusted
mid-execution~\cite{Geist1997:Cumulvs}.
The idea dates to the early
1990s~\cite{Mille1998:SCIRun,Geist1997:Cumulvs} and continues through modern in-situ visualization
tools~\cite{Ayach2015:Catalyst,Larse2022:Ascent} --- mostly fixed-function pipelines over
domain-specific data models.
% SCIRun, RealityGrid
Many systems recognized that steering distributed simulations
requires consistency guarantees: CUMULVS~\cite{Geist1997:Cumulvs} and RMOST~\cite{Loren2008:RMOST}
observe that synchronization boundaries are
the natural points at which updates may be applied, with RMOST formalizing this as a
DSM-based (\emph{distributed shared memory})
consistency model.
Pathfinder~\cite{Hart1999:Pathfinder} proposes coordinated steering via
transaction-based consistency detection with rollback.
However, none
addresses the full problem: a scalable protocol for consistent control inputs from asynchronous
agents, runtime-programmable analytics, and standard data models for the ad-hoc aggregations that
observability
requires.

\paragraph{Performance Analysis Infrastructure}
MRNet~\cite{Roth2003:MRNet} implemented a TBON (\emph{tree-based overlay
  network})~\cite{Mille2008:TBONs} for scalable tool communication, used as a reduction
backend for commercial debuggers~\cite{:TotalView}, but its compiled packet-based filter model
was not designed for analytics.
Chimbuko~\cite{Kelly2020:Chimbuko} implements in-situ trace analysis over TAU~\cite{Shend2006:TAU}
and ADIOS2~\cite{Godoy2020:ADIOS2}, but does not preserve cross-rank
causality, provides no consistent control path, and its analytics are
fixed-function.
Trace analytics increasingly use dataframe-style
interfaces~\cite{Boehm2016:Caliper,Bhate2019:Hatchet,Bhate2024:Pipit,Devar2024:DFTracer},
but require expensive ETL from per-rank logs and operate exclusively
post-hoc~\cite{Yildi2023:IOMax,Jain2025:AMR}.

\paragraph{ML Training Observability}
The ML training community has independently converged on the same problem.
As clusters scale to
tens of thousands of GPUs, the same BSP fragility has driven a wave of recent systems for
detecting and localizing training performance
anomalies~\cite{Wu2025:GREYHOUND,Yao2025:Holmes,Cui2026:Flare,Chen2025:Reveal,Deng2025:Mycroft},
diagnosing faulty hardware~\cite{Deng2025:Minder,Dong2025:Aegis}, and enforcing training
correctness~\cite{Jiang2025:TrainCheck}.
Each builds bespoke infrastructure for subsets of collection, detection, and actuation
from scratch, using selective tracing and event streaming~\cite{Kreps2011:Kafka} to limit telemetry
volume~\cite{Deng2025:Mycroft}\,
or relies on heuristics to trigger collection of local counters if pathologies are
detected~\cite{Wu2025:GREYHOUND,Dong2025:Aegis}.

% Computational steering has traditionally been used to describe classes of scientific simulations
% where the simulation output from a running application is observed, and parameters are
% interactively adjusted.~\cite{Geist1997:Cumulvs}

% Pathfinder provides algorithms, recognizes the need for consistent steering, but uses a central
% coordinator that detects consistency and intervenes to rollback if inconsistency is
% detected~\cite{Hart1999:Pathfinder}. Central coordinators and rollbacks add scalability and
% complexity challenges. Coordinator implements rollbacks.

% Many systems have been proposed, mostly visualization and paramter-oriented.

% RMOST~\cite{Loren2008:RMOST} proposes a distributed shared memory (DSM) abstraction to enable
% steering of simulations Defines consistency models in terms of this common DSM. Recognizing that
% steering inputs should only become visible at synchornization points. Scalability unclear, DSM
% creates extra synchronization. Agents in the loop are intrinsically asynchronous and scalability
% is important.

% In-situ visualization tools widely used. Fixed-function pipelines. Custom data models. Views of
% domain progress. Telemetry requires arbitrary ad-hoc aggregations. Hard to do.

% Except CumulvsRMOST Systems like Cumulvs and RMOST recognize that synchronization defines the
% natural points at which updates may be applied and current state may be read. But neither proposes
% a scalable implementation that model asynchronous agents, or provide standard data models or
% runtime programmable aggregations.

% MRNet is a tree-based overlay with fixed compile filters.

% Ascent uses a dataflow model where mesh data is transformed into scenes to make views and extracts
% to persist data. But foramts like VTK are optimized for meshes and not for general purpose
% analytics. Chimbuko~\cite{Kelly2020:Chimbuko} implements an in-situ trace analysis using TAU and
% ADIOS2. Does not preserve causality, consistent steering, fixed function analytics etc.

% Many systems implement dataframe-based trace
% analytics~\cite{Boehm2016:Caliper,Bhate2019:Hatchet,Bhate2024:Pipit,Devar2024:DFTracer} but
% require expensive ETL to convert traces into these analytical
% formats~\cite{Yildi2023:IOMax,Jain2025:AMR}.
